
% ============================================================
% Tabel 2: Pemetaan Masalah Pengguna
% ============================================================
\begin{table}[h!]
\centering
\caption{Pemetaan masalah pengguna ke kebutuhan solusi \emph{GraphRAG Framework}.}
\label{tab:kebutuhan-pengguna}
\begin{tabularx}{\textwidth}{|c|>{\raggedright\arraybackslash}X|>{\raggedright\arraybackslash}X|>{\raggedright\arraybackslash}X|}
\hline
\textbf{No.} & \textbf{Masalah Pengguna} & \textbf{Akar Penyebab (Gap)} &
\textbf{Kebutuhan Solusi} \\
\hline
M1 & LLM memberikan jawaban keliru atau halusinasi pada domain Framework
    \& Library &
    G1: Absennya \emph{grounding} eksternal &
    Jawaban di-\emph{ground} pada konten SO terverifikasi \\
\hline
M2 & Kualitas jawaban tidak mencerminkan konsensus komunitas SO &
    G2: Sinyal kepercayaan tidak dieksploitasi &
    Bobot kepercayaan (\emph{upvote}, \emph{accepted}) dikodekan dalam KG \\
\hline
M3 & Tidak dapat menelusuri dari mana jawaban berasal &
    G1 + G2: Tidak ada mekanisme \emph{citation} &
    Setiap jawaban menyertakan referensi ke SO thread \\
\hline
M4 & Jawaban tidak dapat menangani pertanyaan relasional kompleks &
    G3: Tidak ada penalaran berbasis struktur graf &
    Traversal graf untuk menghubungkan entitas terkait \\
\hline
M5 & Pengetahuan LLM tidak mencakup versi \emph{library} terbaru &
    G1 + G2: \emph{Knowledge cutoff} tanpa mekanisme pembaruan &
    Pembaruan KG tanpa \emph{re-training} LLM \\
\hline
\end{tabularx}
\end{table}